\subsection{集合乘积的结合性}

命题7。设 \((\mathbf{X}_t)_{t \in \mathbf{I}}\) 设为一个集合族，其指标集 I 不是空集。设 \((\mathrm{J}_{\lambda})_{\lambda \in \mathbf{L}}\) 是 I 的一个划分。则映射

\[
f \rightarrow (\mathrm{pr} _ {\mathbf {J} _ {\lambda}} f) _ {\lambda \in \mathbf {L}}
\]

的 \(\prod_{i\in I}X_i\) 到乘积集合 \(\prod_{\lambda \in L}\left(\prod_{i\in J_\lambda}X_i\right)\) 是双射（集合乘积的“结合性”）。

由 \(\mathbf{pr}_{\mathbf{J}_{\lambda}}f\) 作为图形为 \(f\) 的函数的限制的图形这一解释（第 4 节），可知映射 \(f \to (\mathbf{pr}_{\mathbf{J}_{\lambda}}f)_{\lambda \in \mathbf{L}}\) 是双射这一陈述意味着，对于每个族 \((v_{\lambda})_{\lambda \in \mathbf{L}}\) ，其中 \(v_{\lambda}\) 是从 \(\mathbf{J}_{\lambda}\) 到 \(\bigcup_{i\in I}\mathbf{X}_i\) 的映射，存在唯一的映射 \(u\)

将 I 映入 \(\bigcup_{i\in I}X_i\) ，使得 \(v_{\lambda}\) 是 \(u\) 到 \(J_{\lambda}\) 的限制，对于每一个 \(\lambda\in L\) 。但这是假设 \((J_{\lambda})_{\lambda\in L}\) 是 I 的一个划分（§4，命题 8）的推论。

命题 7 中定义的双射及其逆双射称为典范双射。

\minorheading{注记}

\begin{enumerate}
\def\labelenumi{(\arabic{enumi})}
\tightlist
\item
  让 \(\alpha, \beta\) 是两个不同的对象，令 \((J_{\lambda})_{\lambda \in \{\alpha, \beta\}}\) 是I的一个划分为两个集合 \(J_{\alpha}, J_{\beta}\) 因此我们得到一个一一映射（仍称为典范映射）的乘积 \(\prod_{i \in I} X_i\) 到上 \(\left(\prod_{i \in J_{\alpha}} X_i\right) \times \left(\prod_{i \in J_{\beta}} X_i\right)\) 通过构成典范映射的复合 \(\prod_{\lambda \in \{\alpha, \beta\}} \left(\prod_{i \in J_{\lambda}} X_i\right)\) 到上 \(\left(\prod_{i \in J_{\alpha}} X_i\right) \times \left(\prod_{i \in J_{\beta}} X_i\right)\) 以及 \(\prod_{i \in I} X_i\) 到上 \(\prod_{\lambda \in \{\alpha, \beta\}} \left(\prod_{i \in J_{\lambda}} X_i\right)\) 。如果 \(X_i\) 是一个集合，对于每个 \(i \in J_{\beta}\) ，那么 \(\operatorname{pr}_{J_{\alpha}}\) 是 \(\prod_{i \in I} X_i\) 到上 \(\prod_{i \in J_{\beta}} X_i\) .
\end{enumerate}

\[
\iota \in J _ {\alpha}
\]

\begin{enumerate}
\def\labelenumi{(\arabic{enumi})}
\setcounter{enumi}{1}
\tightlist
\item
  让 \(\alpha, \beta, \gamma\) 是三个对象，其中任意两个都不相等（存在这样的三个对象；例如， \(\phi, \{\phi\}, \{\{\phi'\}\}\) ），并设 A、B、C 为集合。考虑函数图 \(\{(\alpha, A), (\beta, B), (\gamma, C)\}\) ，即集合族 \((X_{t})_{t \in \{\alpha, \beta, \gamma\}}\) 使得 \(X_{\alpha} = A\) , \(X_{\beta} = B\) , \(X_{\gamma} = C\) 。对于由两个集合 \(\{\alpha, \beta, \gamma\}\) 构成的 \(\{\alpha, \beta\}\) 并且 \(\{\gamma\}\) 的分划，存在一个从 \(\prod X_{t}\) 到乘积的典范双射
\end{enumerate}

\(\iota \in \{\alpha, \beta, \gamma\}\)

\[
\left(\prod_ {\iota \in \{\alpha , \beta \}} \mathbf {X} _ {\iota}\right) \times \mathbf {X} _ {\gamma} ^ {\{\gamma \}},
\]

因此是一个双射（同样称为典范的） \(\prod_{\iota \in \{\alpha, \beta, \gamma\}} \mathbf{X}_{\iota}\) 到上 \(\mathbf{A} \times \mathbf{B} \times \mathbf{C}\) ( \(\S 2\) ，编号2）它将每个图 \(f \in \prod_{\iota \in \{\alpha, \beta, \gamma\}} \mathbf{X}_{\iota}\) 变换为元素 \((f(\alpha), f(\beta), f(\gamma))\) 的 \(\mathbf{A} \times \mathbf{B} \times \mathbf{C}\) 。根据命题4，因此我们可以将六个集合中的任意两个 \(\mathbf{A} \times \mathbf{B} \times \mathbf{C}\) , \(\mathbf{B} \times \mathbf{C} \times \mathbf{A}\) , \(\mathbf{C} \times \mathbf{A} \times \mathbf{B}\) , \(\mathbf{B} \times \mathbf{A} \times \mathbf{C}\) , \(\mathbf{A} \times \mathbf{C} \times \mathbf{B}\) , \(\mathbf{C} \times \mathbf{B} \times \mathbf{A}\) 建立一一对应关系。

